\documentclass[11pt, a4paper]{article}
\usepackage[a4paper, top=2.5cm, bottom=2.5cm, left=2cm, right=2cm]{geometry}
\usepackage{fontspec}
\usepackage[english, bidi=basic, provide=*]{babel}
\babelprovide[import, onchar=ids fonts]{english}
\babelfont{rm}{TeX Gyre Termes}
\babelfont{sf}{TeX Gyre Heros}
\babelfont{tt}{TeX Gyre Cursor}
\usepackage{enumitem}
\usepackage{amsmath}

\title{\textbf{Chapter 4: Carbon and the Molecular Diversity of Life \\ 75 Multiple Choice Questions}}
\author{}
\date{}

\begin{document}

\maketitle

\begin{enumerate}

% ---------------------------------------
% 4.1 Organic Chemistry & Origin of Life (1-10)
% ---------------------------------------

\item \textbf{The study of carbon compounds is known as:}
\begin{enumerate}[label=\Alph*.]
    \item Inorganic chemistry
    \item Biochemistry
    \item Organic chemistry
    \item Physical chemistry
    \item Analytical chemistry
\end{enumerate}
\textbf{Answer:} C \\
\textbf{Explanation:} Organic chemistry is specifically defined as the study of carbon compounds, regardless of whether they are produced naturally or synthetically.

\item \textbf{Which element is the foundation of all known life due to its ability to form large, complex, and diverse molecules?}
\begin{enumerate}[label=\Alph*.]
    \item Oxygen
    \item Nitrogen
    \item Phosphorus
    \item Carbon
    \item Hydrogen
\end{enumerate}
\textbf{Answer:} D \\
\textbf{Explanation:} Carbon's tetravalence allows it to act as an intersection point from which a molecule can branch off in up to four directions.

\item \textbf{The historical belief that physical and chemical laws do NOT apply to living things, and that organic compounds only arise in organisms, was known as:}
\begin{enumerate}[label=\Alph*.]
    \item Mechanism
    \item Vitalism
    \item Reductionism
    \item Emergent theory
    \item Natural selection
\end{enumerate}
\textbf{Answer:} B \\
\textbf{Explanation:} Vitalism was the idea that a "life force" outside the jurisdiction of physical and chemical laws was responsible for producing organic compounds.

\item \textbf{Which scientific view eventually shifted the focus away from vitalism by asserting that physical and chemical laws govern all natural phenomena?}
\begin{enumerate}[label=\Alph*.]
    \item Mechanism
    \item Creationism
    \item Isomerism
    \item Darwinism
    \item Entropy
\end{enumerate}
\textbf{Answer:} A \\
\textbf{Explanation:} Mechanism is the view that all natural phenomena, including the processes of life, are governed by the same physical and chemical laws.

\item \textbf{In 1953, Stanley Miller's classic experiment demonstrated that:}
\begin{enumerate}[label=\Alph*.]
    \item Life originated from outer space.
    \item Carbon is the only element capable of forming double bonds.
    \item Organic molecules could be synthesized abiotically under conditions mimicking early Earth.
    \item DNA is the genetic material.
    \item Vitalism is the correct theory of life.
\end{enumerate}
\textbf{Answer:} C \\
\textbf{Explanation:} Miller's experiment produced amino acids and other complex organic compounds spontaneously from simple inorganic gases ($H_2$, $NH_3$, $CH_4$, $H_2O$) and an electrical spark.

\item \textbf{In Miller's experiment, the electrical sparks served as a stand-in for what natural phenomenon on early Earth?}
\begin{enumerate}[label=\Alph*.]
    \item Volcanoes
    \item Tides
    \item Lightning
    \item Sunlight
    \item Geothermal vents
\end{enumerate}
\textbf{Answer:} C \\
\textbf{Explanation:} The sparks provided the intense energy needed to drive the endergonic chemical reactions, mimicking the frequent lightning storms on early Earth.

\item \textbf{Which scientist is credited with accidentally synthesizing urea in the laboratory in 1828, providing the first major blow to vitalism?}
\begin{enumerate}[label=\Alph*.]
    \item Stanley Miller
    \item Friedrich Wöhler
    \item Charles Darwin
    \item Hermann Kolbe
    \item Jöns Jakob Berzelius
\end{enumerate}
\textbf{Answer:} B \\
\textbf{Explanation:} Wöhler tried to make an inorganic salt (ammonium cyanate) but inadvertently produced urea, an organic compound found in animal urine.

\item \textbf{What was the significance of Hermann Kolbe's synthesis of acetic acid?}
\begin{enumerate}[label=\Alph*.]
    \item It proved that water is essential for life.
    \item It showed that organic molecules could be formed from pure elements, further disproving vitalism.
    \item It created the first synthetic amino acid.
    \item It proved that life could not have originated on Earth.
    \item It was the first use of an enzyme in a laboratory.
\end{enumerate}
\textbf{Answer:} B \\
\textbf{Explanation:} Kolbe made acetic acid from pure inorganic elements, proving definitively that a "life force" was not required to assemble organic molecules.

\item \textbf{Which of the following elements is most frequently the bonding partner of carbon in organic molecules?}
\begin{enumerate}[label=\Alph*.]
    \item Sodium
    \item Iron
    \item Hydrogen
    \item Calcium
    \item Chlorine
\end{enumerate}
\textbf{Answer:} C \\
\textbf{Explanation:} Carbon skeletons are most commonly saturated with hydrogen atoms, forming hydrocarbons.

\item \textbf{Recent analyses of meteorites (like the Murchison meteorite) have revealed the presence of:}
\begin{enumerate}[label=\Alph*.]
    \item Living bacteria
    \item Fossilized insects
    \item Pure diamond cores
    \item Amino acids and other organic compounds
    \item High levels of vitalistic energy
\end{enumerate}
\textbf{Answer:} D \\
\textbf{Explanation:} The presence of complex organic molecules in meteorites suggests that the chemical processes producing them are not unique to Earth.

% ---------------------------------------
% 4.2 Carbon Bonding & Skeletons (11-40)
% ---------------------------------------

\item \textbf{Carbon has an atomic number of 6. How many valence electrons does it have?}
\begin{enumerate}[label=\Alph*.]
    \item 2
    \item 4
    \item 6
    \item 8
    \item 12
\end{enumerate}
\textbf{Answer:} B \\
\textbf{Explanation:} A carbon atom has 2 electrons in its inner shell and 4 electrons in its outer (valence) shell.

\item \textbf{Because carbon has 4 valence electrons, it completes its valence shell by forming a maximum of how many covalent bonds?}
\begin{enumerate}[label=\Alph*.]
    \item 1
    \item 2
    \item 3
    \item 4
    \item 8
\end{enumerate}
\textbf{Answer:} D \\
\textbf{Explanation:} Carbon acts as an intersection point from which a molecule can branch off in four directions, known as tetravalence.

\item \textbf{When a carbon atom forms four single covalent bonds, the bonds angle toward the corners of an imaginary:}
\begin{enumerate}[label=\Alph*.]
    \item Cube
    \item Octahedron
    \item Tetrahedron
    \item Sphere
    \item Flat square
\end{enumerate}
\textbf{Answer:} C \\
\textbf{Explanation:} The repulsion between the four electron pairs in the hybrid orbitals forces them to the corners of a tetrahedron, with bond angles of roughly $109.5^\circ$.

\item \textbf{When two carbon atoms are joined by a double bond, how many electron pairs are shared between them?}
\begin{enumerate}[label=\Alph*.]
    \item 1
    \item 2
    \item 3
    \item 4
    \item 6
\end{enumerate}
\textbf{Answer:} B \\
\textbf{Explanation:} A single bond is the sharing of one pair of electrons; a double bond is the sharing of two pairs (4 electrons total).

\item \textbf{What is the resulting 3D shape of a molecule when two carbon atoms are joined by a double bond (e.g., Ethene, $C_2H_4$)?}
\begin{enumerate}[label=\Alph*.]
    \item It forms two linked tetrahedrons.
    \item It forms a straight, linear line.
    \item All atoms attached to those carbons lie in the same flat plane.
    \item It forms a spherical shell.
    \item It forms a rigid cube.
\end{enumerate}
\textbf{Answer:} C \\
\textbf{Explanation:} The double bond restricts the geometry, making the two carbons and the four atoms attached to them co-planar (flat).

\item \textbf{The valences of the four major atomic components of organic molecules—Hydrogen, Oxygen, Nitrogen, and Carbon—are respectively:}
\begin{enumerate}[label=\Alph*.]
    \item 1, 2, 3, 4
    \item 1, 3, 2, 4
    \item 4, 3, 2, 1
    \item 2, 4, 3, 1
    \item 1, 4, 2, 3
\end{enumerate}
\textbf{Answer:} A \\
\textbf{Explanation:} Hydrogen forms 1 bond, Oxygen 2, Nitrogen 3, and Carbon 4. This is the "building code" of molecular architecture.

\item \textbf{In a molecule of carbon dioxide ($CO_2$), the central carbon atom forms how many double bonds?}
\begin{enumerate}[label=\Alph*.]
    \item 0
    \item 1
    \item 2
    \item 3
    \item 4
\end{enumerate}
\textbf{Answer:} C \\
\textbf{Explanation:} Carbon forms a double bond with each of the two oxygen atoms ($O=C=O$), satisfying its valence of 4 and rendering the molecule linear.

\item \textbf{Molecules that consist ONLY of carbon and hydrogen are called:}
\begin{enumerate}[label=\Alph*.]
    \item Carbohydrates
    \item Hydrocarbons
    \item Amines
    \item Alcohols
    \item Polymers
\end{enumerate}
\textbf{Answer:} B \\
\textbf{Explanation:} Hydrocarbons, like petroleum components and methane, contain exclusively carbon and hydrogen atoms.

\item \textbf{Which of the following is a key chemical property of hydrocarbons?}
\begin{enumerate}[label=\Alph*.]
    \item They are highly hydrophilic.
    \item They dissolve easily in water.
    \item They form extensive hydrogen bonds.
    \item They are nonpolar and hydrophobic.
    \item They are highly acidic.
\end{enumerate}
\textbf{Answer:} D \\
\textbf{Explanation:} The C-H bond is nonpolar, making hydrocarbon chains hydrophobic (water-repelling).

\item \textbf{Biological molecules like fats have long hydrocarbon tails. What is the primary biological function of these hydrocarbon tails?}
\begin{enumerate}[label=\Alph*.]
    \item To dissolve in the bloodstream
    \item To act as enzymes
    \item To store large amounts of energy
    \item To carry genetic information
    \item To buffer pH
\end{enumerate}
\textbf{Answer:} C \\
\textbf{Explanation:} Hydrocarbon bonds contain a large amount of potential energy. The breakdown of fat releases significantly more energy than carbohydrates.

\item \textbf{Carbon skeletons can vary in all of the following ways EXCEPT:}
\begin{enumerate}[label=\Alph*.]
    \item Length
    \item Branching
    \item Presence of ring structures
    \item Position of double bonds
    \item The type of covalent bonds connecting purely carbon atoms (they can only form single bonds)
\end{enumerate}
\textbf{Answer:} E \\
\textbf{Explanation:} Carbon atoms can easily form single, double, and even triple bonds with other carbon atoms. The other options are valid ways carbon skeletons vary.

\item \textbf{Which of the following molecules represents a straight-chain (unbranched) hydrocarbon?}
\begin{enumerate}[label=\Alph*.]
    \item Isobutane
    \item Benzene
    \item Propane
    \item Cyclohexane
    \item Cholesterol
\end{enumerate}
\textbf{Answer:} C \\
\textbf{Explanation:} Propane ($C_3H_8$) is a simple, unbranched three-carbon chain.

\item \textbf{Which of the following molecules represents a hydrocarbon arranged in a ring?}
\begin{enumerate}[label=\Alph*.]
    \item Butane
    \item Ethene
    \item Cyclohexane
    \item Isobutane
    \item 1-Butene
\end{enumerate}
\textbf{Answer:} C \\
\textbf{Explanation:} Cyclohexane is a six-carbon ring, representing variation in carbon skeleton structure.

\item \textbf{What determines whether a carbon atom's covalent bonds will have a tetrahedral or a planar geometry?}
\begin{enumerate}[label=\Alph*.]
    \item The presence or absence of double bonds
    \item The temperature of the environment
    \item The number of neutrons in the carbon nucleus
    \item Whether it bonds to oxygen or nitrogen
    \item The pH of the solution
\end{enumerate}
\textbf{Answer:} A \\
\textbf{Explanation:} Single bonds result in sp3 hybridization (tetrahedral), whereas a double bond results in sp2 hybridization (planar geometry).

\item \textbf{Why are hydrocarbons typically poor solvents for polar substances like salt?}
\begin{enumerate}[label=\Alph*.]
    \item Because they are highly electronegative.
    \item Because they lack partial charges to form hydration shells.
    \item Because they are too heavy.
    \item Because they form strong ionic bonds with salt.
    \item Because they are exclusively found as gases.
\end{enumerate}
\textbf{Answer:} B \\
\textbf{Explanation:} Hydrocarbons are nonpolar. Without partial positive or negative charges, they cannot surround and separate ions or polar molecules.

% ---------------------------------------
% Isomers (26-40)
% ---------------------------------------

\item \textbf{Compounds that have the same molecular formula but different structures and properties are called:}
\begin{enumerate}[label=\Alph*.]
    \item Isotopes
    \item Polymers
    \item Isomers
    \item Functional groups
    \item Enantiomers only
\end{enumerate}
\textbf{Answer:} C \\
\textbf{Explanation:} Isomers have identical formulas (same number of atoms) but different architectural arrangements.

\item \textbf{Pentane and isopentane both have the formula $C_5H_{12}$, but isopentane is branched while pentane is a straight chain. These are examples of:}
\begin{enumerate}[label=\Alph*.]
    \item Enantiomers
    \item Cis-trans isomers
    \item Geometric isomers
    \item Structural isomers
    \item Radioactive isotopes
\end{enumerate}
\textbf{Answer:} D \\
\textbf{Explanation:} Structural isomers differ in the covalent arrangement of their atoms (e.g., branching).

\item \textbf{As the carbon skeleton increases in size, the number of possible structural isomers:}
\begin{enumerate}[label=\Alph*.]
    \item Decreases
    \item Remains constant
    \item Increases dramatically
    \item Approaches zero
    \item Becomes exactly 3
\end{enumerate}
\textbf{Answer:} C \\
\textbf{Explanation:} With more atoms to arrange, the mathematical possibilities for branching and atom placement increase exponentially.

\item \textbf{Cis-trans isomers (geometric isomers) require the presence of which structural feature?}
\begin{enumerate}[label=\Alph*.]
    \item A carbon-carbon single bond
    \item An asymmetric carbon atom
    \item A hydroxyl group
    \item A carbon-carbon double bond
    \item A nitrogen atom
\end{enumerate}
\textbf{Answer:} D \\
\textbf{Explanation:} Double bonds are inflexible and do not allow free rotation, locking the attached atoms into specific spatial arrangements relative to the double bond.

\item \textbf{In a cis isomer, the two designated atoms or groups (like two Xs) are located:}
\begin{enumerate}[label=\Alph*.]
    \item On opposite sides of the double bond.
    \item On the same side of the double bond.
    \item On the same carbon atom.
    \item In a separate molecule.
    \item Above and below a central ring.
\end{enumerate}
\textbf{Answer:} B \\
\textbf{Explanation:} "Cis" means on the same side. "Trans" means across or on opposite sides.

\item \textbf{Why don't single covalent bonds between carbon atoms form cis-trans isomers?}
\begin{enumerate}[label=\Alph*.]
    \item Because they are too weak.
    \item Because they allow for free rotation of the atoms around the bond axis.
    \item Because they are planar.
    \item Because single bonds only occur in rings.
    \item Because they are always asymmetrical.
\end{enumerate}
\textbf{Answer:} B \\
\textbf{Explanation:} Single bonds allow the bonded groups to rotate freely, so "same side" vs "opposite side" configurations quickly interchange and are not fixed.

\item \textbf{Enantiomers are isomers that:}
\begin{enumerate}[label=\Alph*.]
    \item Differ in the arrangement of covalent bonds.
    \item Are mirror images of each other.
    \item Differ in the placement of a double bond.
    \item Have different molecular formulas.
    \item Always possess exactly identical biological functions.
\end{enumerate}
\textbf{Answer:} B \\
\textbf{Explanation:} Enantiomers are left-handed and right-handed versions of a molecule, acting as non-superimposable mirror images.

\item \textbf{For a carbon atom to be considered "asymmetric" and thus capable of forming enantiomers, it must be bonded to:}
\begin{enumerate}[label=\Alph*.]
    \item Four identical atoms.
    \item A double bond and two single bonds.
    \item Four different atoms or groups of atoms.
    \item Only hydrogen atoms.
    \item At least one oxygen atom.
\end{enumerate}
\textbf{Answer:} C \\
\textbf{Explanation:} An asymmetric (chiral) carbon must have four distinct groups attached to it to lack a plane of symmetry.

\item \textbf{In pharmacology, it is common that only one enantiomer of a drug is active. Why?}
\begin{enumerate}[label=\Alph*.]
    \item One enantiomer usually degrades into a gas.
    \item The inactive enantiomer lacks carbon atoms.
    \item Only one enantiomer has the correct spatial shape to bind to the specific receptor molecule in the body.
    \item The body temperature denatures one of the enantiomers.
    \item One enantiomer is typically a strong acid, while the other is a strong base.
\end{enumerate}
\textbf{Answer:} C \\
\textbf{Explanation:} Biological receptors are highly specific 3D structures. Like a right glove only fitting a right hand, receptors usually only bind one enantiomer.

\item \textbf{L-Dopa is an effective treatment for Parkinson's disease, whereas D-Dopa has no effect. This illustrates the importance of:}
\begin{enumerate}[label=\Alph*.]
    \item Structural isomers
    \item Cis-trans isomers
    \item Enantiomers
    \item Isotopic decay
    \item Hydrophobic interactions
\end{enumerate}
\textbf{Answer:} C \\
\textbf{Explanation:} L-Dopa and D-Dopa are enantiomers. The enzymes/receptors in the brain only recognize the L form.

\item \textbf{Methamphetamine occurs as two enantiomers. One is a highly addictive illegal drug (crank), while the other is:}
\begin{enumerate}[label=\Alph*.]
    \item A deadly poison
    \item An over-the-counter vapor inhaler for nasal congestion
    \item A treatment for Parkinson's disease
    \item A common food preservative
    \item A vital vitamin
\end{enumerate}
\textbf{Answer:} B \\
\textbf{Explanation:} The active enantiomer causes intense neurological effects, while the mirror image has completely different, much weaker physiological effects (sold as an inhaler).

\item \textbf{Thalidomide was prescribed in the 1950s to treat morning sickness. It was later discovered to exist as a mixture of enantiomers, where one caused:}
\begin{enumerate}[label=\Alph*.]
    \item Extreme hair loss
    \item Blindness in the mother
    \item Severe birth defects in embryos
    \item Spontaneous combustion
    \item Accelerated aging
\end{enumerate}
\textbf{Answer:} C \\
\textbf{Explanation:} Thalidomide's tragic history is a primary example of how differently enantiomers can behave biologically.

\item \textbf{If a molecule has three asymmetric carbons, what is the maximum number of possible stereoisomers (enantiomers + diastereomers)?}
\begin{enumerate}[label=\Alph*.]
    \item 3
    \item 6
    \item 8
    \item 9
    \item 12
\end{enumerate}
\textbf{Answer:} C \\
\textbf{Explanation:} The rule for maximum stereoisomers is $2^n$, where n is the number of asymmetric carbons. $2^3 = 8$.

\item \textbf{Which property is identical between two enantiomers of a given molecule?}
\begin{enumerate}[label=\Alph*.]
    \item Biological activity
    \item Smell and taste
    \item Interaction with polarized light
    \item Chemical formula and molecular weight
    \item Ability to bind to an enzyme
\end{enumerate}
\textbf{Answer:} D \\
\textbf{Explanation:} Enantiomers are chemically identical in formula, weight, and general physical properties, but differ in 3D interactions (like receptor binding and rotating plane-polarized light).

\item \textbf{Ibuprofen is a drug sold as a mixture of enantiomers. Which of the following is true regarding Ibuprofen?}
\begin{enumerate}[label=\Alph*.]
    \item Both enantiomers are equally effective at reducing inflammation.
    \item The S-enantiomer is effectively 100 times more active than the R-enantiomer.
    \item The R-enantiomer is toxic.
    \item The drug is purely a cis-trans isomer mixture.
    \item It contains no asymmetric carbons.
\end{enumerate}
\textbf{Answer:} B \\
\textbf{Explanation:} The S-enantiomer is the biologically active form that inhibits the enzymes responsible for inflammation.

% ---------------------------------------
% 4.3 Functional Groups (41-71)
% ---------------------------------------

\item \textbf{Chemical groups attached to a carbon skeleton that participate directly in chemical reactions are known as:}
\begin{enumerate}[label=\Alph*.]
    \item Isotopes
    \item Inert groups
    \item Functional groups
    \item Hydrocarbons
    \item Noble gases
\end{enumerate}
\textbf{Answer:} C \\
\textbf{Explanation:} Functional groups have specific chemical properties (e.g., charge, polarity, acidity) that govern molecular behavior.

\item \textbf{Which chemical group is NOT strictly considered a reactive functional group, but rather serves as a recognizable biological tag on molecules like DNA?}
\begin{enumerate}[label=\Alph*.]
    \item Amino group
    \item Hydroxyl group
    \item Methyl group
    \item Sulfhydryl group
    \item Carboxyl group
\end{enumerate}
\textbf{Answer:} C \\
\textbf{Explanation:} The methyl group ($-CH_3$) is nonpolar and unreactive, but modifying molecules with it drastically alters their recognition by enzymes (e.g., DNA methylation affects gene expression).

\item \textbf{The functional group represented by $-OH$ is called the:}
\begin{enumerate}[label=\Alph*.]
    \item Carbonyl group
    \item Carboxyl group
    \item Hydroxyl group
    \item Sulfhydryl group
    \item Amino group
\end{enumerate}
\textbf{Answer:} C \\
\textbf{Explanation:} The hydroxyl group consists of an oxygen bonded to a hydrogen.

\item \textbf{Organic molecules containing hydroxyl groups are broadly classified as:}
\begin{enumerate}[label=\Alph*.]
    \item Thiols
    \item Alcohols
    \item Amines
    \item Ketones
    \item Carboxylic acids
\end{enumerate}
\textbf{Answer:} B \\
\textbf{Explanation:} The presence of an $-OH$ group defines an alcohol (e.g., ethanol).

\item \textbf{Which property is characteristic of a hydroxyl group?}
\begin{enumerate}[label=\Alph*.]
    \item It acts as a strong base.
    \item It makes the molecule highly hydrophobic.
    \item It is polar and can form hydrogen bonds with water.
    \item It stabilizes protein structure via cross-linking.
    \item It donates $H^+$ ions to lower pH.
\end{enumerate}
\textbf{Answer:} C \\
\textbf{Explanation:} The electronegative oxygen draws electrons, creating a polar group that makes alcohols soluble in water.

\item \textbf{A carbonyl group ($>C=O$) located completely within a carbon skeleton (not at the end) characterizes a class of compounds called:}
\begin{enumerate}[label=\Alph*.]
    \item Aldehydes
    \item Ketones
    \item Thiols
    \item Organic phosphates
    \item Amines
\end{enumerate}
\textbf{Answer:} B \\
\textbf{Explanation:} If the carbonyl is internal, the compound is a ketone (e.g., acetone). If it is at the end of the skeleton, it is an aldehyde.

\item \textbf{Sugars that possess an aldehyde group are termed:}
\begin{enumerate}[label=\Alph*.]
    \item Ketoses
    \item Amines
    \item Aldoses
    \item Carboxylics
    \item Thiols
\end{enumerate}
\textbf{Answer:} C \\
\textbf{Explanation:} Sugars are categorized by the position of their carbonyl group into aldoses (terminal) and ketoses (internal).

\item \textbf{The functional group $-COOH$ is known as the:}
\begin{enumerate}[label=\Alph*.]
    \item Carbonyl group
    \item Hydroxyl group
    \item Carboxyl group
    \item Phosphate group
    \item Methyl group
\end{enumerate}
\textbf{Answer:} C \\
\textbf{Explanation:} The carboxyl group combines a carbonyl ($C=O$) and a hydroxyl ($-OH$) on the same carbon.

\item \textbf{Why does the carboxyl group act as an acid?}
\begin{enumerate}[label=\Alph*.]
    \item It readily absorbs $H^+$ from solution.
    \item It completely nonpolarizes the molecule.
    \item The covalent bond between oxygen and hydrogen is so polar that $H^+$ easily dissociates.
    \item It generates $OH^-$ ions in water.
    \item It destroys buffers.
\end{enumerate}
\textbf{Answer:} C \\
\textbf{Explanation:} The highly electronegative oxygen atoms pull electrons away from hydrogen, allowing the hydrogen ion to leave easily, donating a proton to the solution.

\item \textbf{A compound containing a carboxyl group is known as a(n):}
\begin{enumerate}[label=\Alph*.]
    \item Ketone
    \item Alcohol
    \item Thiol
    \item Carboxylic acid (organic acid)
    \item Amine
\end{enumerate}
\textbf{Answer:} D \\
\textbf{Explanation:} Acetic acid (vinegar) is a classic example of a carboxylic acid.

\item \textbf{The amino group consists of:}
\begin{enumerate}[label=\Alph*.]
    \item A carbon double-bonded to an oxygen.
    \item An oxygen bonded to a hydrogen.
    \item A nitrogen bonded to two hydrogens.
    \item A sulfur bonded to a hydrogen.
    \item A phosphorus surrounded by oxygens.
\end{enumerate}
\textbf{Answer:} C \\
\textbf{Explanation:} The amino group is written as $-NH_2$.

\item \textbf{Because the amino group can pick up an $H^+$ from the surrounding solution, it acts as a(n):}
\begin{enumerate}[label=\Alph*.]
    \item Acid
    \item Base
    \item Buffer
    \item Solvent
    \item Catalyst
\end{enumerate}
\textbf{Answer:} B \\
\textbf{Explanation:} By accepting a proton, the amino group reduces the $H^+$ concentration of the solution, which is the definition of a base.

\item \textbf{Amino acids are the building blocks of proteins. They are unique because every amino acid contains at least one amino group and one \_\_\_\_\_\_\_ group.}
\begin{enumerate}[label=\Alph*.]
    \item Sulfhydryl
    \item Phosphate
    \item Methyl
    \item Carboxyl
    \item Carbonyl
\end{enumerate}
\textbf{Answer:} D \\
\textbf{Explanation:} The name "amino acid" refers to the presence of both the basic amino group ($-NH_2$) and the acidic carboxyl group ($-COOH$).

\item \textbf{The sulfhydryl group is represented by the formula:}
\begin{enumerate}[label=\Alph*.]
    \item $-OH$
    \item $-SH$
    \item $-NH_2$
    \item $-CH_3$
    \item $-COOH$
\end{enumerate}
\textbf{Answer:} B \\
\textbf{Explanation:} The sulfhydryl group consists of a sulfur atom bonded to a hydrogen atom.

\item \textbf{Compounds containing sulfhydryl groups are called:}
\begin{enumerate}[label=\Alph*.]
    \item Thiols
    \item Alcohols
    \item Amines
    \item Acids
    \item Ketones
\end{enumerate}
\textbf{Answer:} A \\
\textbf{Explanation:} "Thio-" refers to sulfur; thus, these alcohols-like molecules with sulfur instead of oxygen are called thiols.

\item \textbf{What is a crucial biological function of the sulfhydryl group?}
\begin{enumerate}[label=\Alph*.]
    \item It acts as the primary energy carrier in ATP.
    \item It serves as a strong acid in the stomach.
    \item Two sulfhydryl groups can react to form a cross-link that stabilizes protein structures.
    \item It acts as a biological tag to turn genes off.
    \item It gives sugar its sweet taste.
\end{enumerate}
\textbf{Answer:} C \\
\textbf{Explanation:} Disulfide bridges ($-S-S-$) form between cysteine amino acids, giving shape and stability to proteins (e.g., determining hair curliness).

\item \textbf{Which functional group is most directly involved in energy transfers within a cell?}
\begin{enumerate}[label=\Alph*.]
    \item Hydroxyl
    \item Methyl
    \item Phosphate
    \item Sulfhydryl
    \item Amino
\end{enumerate}
\textbf{Answer:} C \\
\textbf{Explanation:} Phosphate groups ($-OPO_3^{2-}$) store and release tremendous amounts of energy, notably as components of ATP.

\item \textbf{When attached to a molecule, the phosphate group generally contributes what type of charge?}
\begin{enumerate}[label=\Alph*.]
    \item A positive charge of +1
    \item A positive charge of +2
    \item A neutral charge
    \item A negative charge (either 1- or 2-)
    \item A shifting dipole charge
\end{enumerate}
\textbf{Answer:} D \\
\textbf{Explanation:} The oxygen atoms in the phosphate group lose their protons in cellular pH, carrying negative charges.

\item \textbf{The methyl group is represented by:}
\begin{enumerate}[label=\Alph*.]
    \item $-NH_2$
    \item $-CH_3$
    \item $-SH$
    \item $-OH$
    \item $>C=O$
\end{enumerate}
\textbf{Answer:} B \\
\textbf{Explanation:} A carbon bonded to three hydrogens constitutes a methyl group.

\item \textbf{The addition of a methyl group to DNA (DNA methylation) affects the molecule primarily by:}
\begin{enumerate}[label=\Alph*.]
    \item Breaking the DNA strand.
    \item Altering the expression of genes.
    \item Causing immediate cell death.
    \item Transforming the DNA into RNA.
    \item Changing the pH of the nucleus.
\end{enumerate}
\textbf{Answer:} B \\
\textbf{Explanation:} Methyl groups act as "tags" that enzymes recognize, often silencing or altering gene expression without reacting chemically.

\item \textbf{The visual and physiological differences between male and female lions are largely due to small differences in the chemical groups attached to a common carbon skeleton found in:}
\begin{enumerate}[label=\Alph*.]
    \item ATP
    \item Amino acids
    \item Sex hormones (estradiol and testosterone)
    \item Carbohydrates
    \item DNA
\end{enumerate}
\textbf{Answer:} C \\
\textbf{Explanation:} Estradiol and testosterone share an identical four-ring carbon skeleton but differ in just a few attached chemical groups (like a methyl vs. hydroxyl group).

% ---------------------------------------
% 4.3 ATP & Review (72-75)
% ---------------------------------------

\item \textbf{ATP stands for:}
\begin{enumerate}[label=\Alph*.]
    \item Adenine thymine phosphate
    \item Adenosine triphosphate
    \item Active transport protein
    \item Adenine triphosphate
    \item Amino threonine phosphate
\end{enumerate}
\textbf{Answer:} B \\
\textbf{Explanation:} ATP consists of an organic molecule called adenosine attached to a string of three phosphate groups.

\item \textbf{When ATP reacts with water, what happens to one of the phosphate groups?}
\begin{enumerate}[label=\Alph*.]
    \item It becomes a methyl group.
    \item It attaches permanently to the adenosine.
    \item It splits off, releasing energy and leaving behind ADP.
    \item It absorbs energy and becomes a fourth phosphate group.
    \item It forms a sulfhydryl cross-link.
\end{enumerate}
\textbf{Answer:} C \\
\textbf{Explanation:} The hydrolysis of ATP breaks the bond of the terminal phosphate group, yielding ADP, inorganic phosphate ($P_i$), and releasing energy used by the cell.

\item \textbf{Which of the following molecules acts as the primary "energy currency" of the cell?}
\begin{enumerate}[label=\Alph*.]
    \item Glucose
    \item DNA
    \item ATP
    \item Cysteine
    \item Methane
\end{enumerate}
\textbf{Answer:} C \\
\textbf{Explanation:} ATP stores the potential to react with water, releasing the specific bursts of energy needed to drive cellular processes.

\item \textbf{If a molecule contains a hydroxyl group, a carbonyl group, and a carbon skeleton of 6 carbons, it is most likely a(n):}
\begin{enumerate}[label=\Alph*.]
    \item Amino acid
    \item Sugar (carbohydrate)
    \item Fat
    \item DNA base
    \item Steroid hormone
\end{enumerate}
\textbf{Answer:} B \\
\textbf{Explanation:} Sugars are characterized by multiple hydroxyl groups and one carbonyl group attached to a carbon skeleton.

\item \textbf{Which two functional groups are always found in amino acids?}
\begin{enumerate}[label=\Alph*.]
    \item Ketone and Aldehyde
    \item Carbonyl and Amino
    \item Carboxyl and Amino
    \item Phosphate and Sulfhydryl
    \item Hydroxyl and Carboxyl
\end{enumerate}
\textbf{Answer:} C \\
\textbf{Explanation:} The name "amino acid" denotes the presence of both an amino group (basic) and a carboxyl group (acidic).

\item \textbf{The flexibility of carbon bonding allows life's molecules to be:}
\begin{enumerate}[label=\Alph*.]
    \item Incredibly diverse in architecture and function.
    \item Highly rigid and unchanging.
    \item Purely linear in all cases.
    \item Composed only of metallic ions.
    \item Exclusively nonpolar.
\end{enumerate}
\textbf{Answer:} A \\
\textbf{Explanation:} Carbon's tetravalence allows for chains, branching, rings, and double bonds, creating immense molecular diversity essential for life.

\end{enumerate}

\end{document}
